\documentclass{article}
\usepackage[UTF8]{ctex}
\usepackage{amsmath,amssymb,amsfonts}
\usepackage{enumitem}
\usepackage{geometry}
\geometry{a4paper, margin=1in}
% === 错题本排版优化（不参与题目标识） ===
\setlength{\parindent}{0pt}
\setlength{\parskip}{0.35em plus 0.10em minus 0.08em}
\setlength{\emergencystretch}{3em}
\setlist{itemsep=0.25em, topsep=0.40em, parsep=0pt}
\allowdisplaybreaks[2]
% === 错题本排版优化结束 ===
\begin{document}
\section*{中山大学2026年高等代数真题}
\section*{144．中山大学2026年高等代数试卷}

\subsection{题目1}
一．（15 分）下面行列式

\[
d(x)=\left|\begin{array}{cccc}
2 x & x & 1 & 2 \\
1 & x & 1 & -1 \\
3 & 2 & x & 1 \\
1 & 1 & 1 & x
\end{array}\right|
\]

中 \(x^{k}\) 的系数记为 \(a_{k}\) ，求 \(a_{2}\) 及 \(a_{3}\) 。

\subsection{题目2}
二．（15 分）已知实数域上关于 \(X=\left(x_{1}, x_{2}, x_{3}\right)^{\mathrm{T}}\) 的线性方程组

\[
\left\{\begin{array}{l}
x_{1}+a x_{2}+x_{3}=4 \\
b x_{1}+x_{2}+x_{3}=3 \\
2 b x_{1}+x_{2}+x_{3}=4
\end{array}\right.
\]

有解，求 \(a, b\) 的值，并求出方程组的解。

\subsection{题目3}
三．（15 分）实矩阵

\[
A=\left(\begin{array}{ccc}
1 & -2 & -4 \\
-2 & x & -2 \\
-4 & -2 & 1
\end{array}\right)
\]

与对角阵 \(\Lambda=\operatorname{diag}\{5,-4, y\}\) 相似，求 \(x, y\) 的值．

\subsection{题目4}
四．（15 分）设 \(A=\left(\begin{array}{lll}0 & 1 & 1 \\ 1 & 0 & 1 \\ 1 & 1 & 0\end{array}\right)\) ，求正交矩阵 \(P\) 及对角矩阵 \(D\) ，使得 \(P^{\mathrm{T}} A P=D\) ．

\subsection{题目5}
五．（15 分）设

\[
A=\left(\begin{array}{ccc}
4 & -15 & 6 \\
1 & -4 & 2 \\
1 & -5 & 3
\end{array}\right) \in M_{3}(\mathbb{C})
\]

求 \(A\) 的 Jordan 标准形 \(J\) 及 \(A\) 的最小多项式 \(g(\lambda)\) ．

\subsection{题目6}
六。（15 分）设 \(A \in M_{n}(\mathbb{C})\) ，任意给定非零列向量 \(\alpha \in \mathbb{C}^{n}\) ，求证：存在次数小 \(n\) 的多项式 \(h(x) \in \mathbb{C}[x]\)，使得 \(h(A) \alpha\) 为 \(A\) 的特征向量。

\subsection{题目7}
七．（15 分）设 \(A, B \in M_{n}(\mathbb{R})\) 为可逆矩阵，满足 \(A B A=B\) 且 \(B A B=A\) ，证明：\(A^{4}=B^{4}=I\) ．

\subsection{题目8}
八．（15 分）设 \(A=\left(a_{i j}\right)_{n \times n} \in M_{n}(\mathbb{R})\) 是正定对称矩阵，其对角元 \(a_{i i}\) 及特征值 \(\lambda_{i}\) 满足

\[
a_{11} \leq a_{22} \leq \cdots \leq a_{n n}, \lambda_{1} \leq \lambda_{2} \leq \cdots \leq \lambda_{n} .
\]

问 \(\lambda_{i} \leq a_{i i}\) 是否对所有 \(1 \leq i \leq n\) 成立？若成立，给出证明；若不成立，举出反例．

\subsection{题目9}
九．（15 分）对非负整数 \(n\) ，已知 \(M_{2 n+1}(\mathbb{R})\) 关于矩阵加法与数乘构成一个实线性空间，设 \(V\) 是\(M_{2 n+1}(\mathbb{R})\) 的一个非零子空间，满足 \(V\) 中任意非零矩阵都可逆。求证： \(\operatorname{dim} V=1\) 。

\subsection{题目10}
十．（15 分）记 \(E_{i j}\) 是 \((i, j)\) 元素是1其余元素是 0 的 \(n\) 阶实矩阵。设 \(\varphi\) 是实列向量空间 \(\mathbb{R}^{n}\) 上的线性变换，定义映射

\[
\begin{gathered}
f_{\varphi}: \mathbb{R}^{n} \rightarrow \mathbb{R}^{n} \\
\alpha \mapsto \sum_{i, j=1}^{n} E_{i j} \varphi\left(E_{j i} \alpha\right)
\end{gathered}
\]

求证：存在常数 \(c_{\varphi}\) ，满足 \(f_{\varphi}=c_{\varphi} \mathrm{Id}, \mathrm{Id}\) 为 \(\mathbb{R}^{n}\) 上的恒等变换。
\end{document}
